% GENERATED FILE. Edit canonical lessons, then run npm run generate:books.
% canonical-source-hash: fnv1a64:3325b3c84cdd22b4
% canonical-lessons: TA-C417-malumi, TA-C417-nankuram, TA-C417-teppam, TA-C417-tattali, TA-C417-karai

\chapter{Sailors and Anchors}
\label{ch:ta-sailors-and-anchors}
\hlchaptermodality{417}

\begin{chapteropening}
\textbf{By the end of this chapter:} \emph{I can talk about sailors, anchors, rafts and the shore, and a boat floundering.}

Complete the last lesson of chapter 417: I can talk about sailors, anchors, rafts and the shore, and a boat floundering.
\end{chapteropening}

\section[\texorpdfstring{\ta{மாலுமி} (mālumi)}{மாலுமி (mālumi)}]{\ta{மாலுமி} (mālumi) — a sailor}

\label{lesson:TA-C417-malumi}



\begin{quote}
Before the new one: say the Tamil for a catamaran, then the Tamil for a net.
\end{quote}

\subsection*{You'll want to know: \ta{மாலுமி}}
\textbf{\ta{மாலுமி}} — \emph{mālumi} — \textquotedblleft{}a sailor\textquotedblright{}.

\begin{quote}
\ta{மாலுமிகள்} \ta{கப்பலைச்} \ta{சுத்தம்} \ta{செய்தார்கள்}.
\end{quote}

\begin{quote}
\emph{mālumikaḷ kappalaic cuttam ceytārkaḷ.}
\end{quote}

\begin{quote}
The sailors cleaned the ship.
\end{quote}

The plural is \ta{மாலுமிகள்}.

\begin{cousinweb}
\textbf{\ta{மாலுமி}} (\emph{mālumi}) is usually traced to Arabic \emph{muʿallim}, \textquotedblleft{}a master, a ship's captain\textquotedblright{}; the path into Tamil is not settled.
\end{cousinweb}

\begin{grammarlens}[title={What you've built}]
One more word for getting around.
\end{grammarlens}

\subsection*{Your turn}
\begin{itemize}
\raggedright
  \item \emph{Say it:} \emph{mālumi}
  \item \emph{Say it:} \emph{mālumi}, once more
  \item \emph{Say it:} \emph{mālumikaḷ kappalaic cuttam ceytārkaḷ.} — The sailors cleaned the ship.
  \item \emph{Recall:} say the Tamil for a catamaran, then the Tamil for a net, then say \emph{mālumi} again
\end{itemize}

\begin{tcolorbox}[breakable,colback=teal!4,colframe=teal!35!black,title={Before you move on}]
What does \ta{மாலுமி} mean? (\textquotedblleft{}A sailor\textquotedblright{}.) Say it once more.
\end{tcolorbox}
\section[\texorpdfstring{\ta{நங்கூரம்} (naṅkūram)}{நங்கூரம் (naṅkūram)}]{\ta{நங்கூரம்} (naṅkūram) — an anchor}

\label{lesson:TA-C417-nankuram}



\begin{quote}
Before the new one: say the Tamil for a net, then the Tamil for a sailor.
\end{quote}

\subsection*{You'll want to know: \ta{நங்கூரம்}}
\textbf{\ta{நங்கூரம்}} — \emph{naṅkūram} — \textquotedblleft{}an anchor\textquotedblright{}.

\begin{quote}
\ta{கப்பல்} \ta{துறைமுகத்துக்கு} \ta{வெளியே} \ta{நங்கூரம்} \ta{போட்டது}.
\end{quote}

\begin{quote}
\emph{kappal tuṟaimukattukku veḷiyē naṅkūram pōṭṭatu.}
\end{quote}

\begin{quote}
The ship dropped anchor outside the harbour.
\end{quote}

\ta{நங்கூரம்} \ta{போடு} is to drop anchor.

\begin{cousinweb}
\textbf{\ta{நங்கூரம்}} (\emph{naṅkūram}) is usually traced to Persian \emph{langar}, \textquotedblleft{}an anchor\textquotedblright{}.
\end{cousinweb}

\begin{grammarlens}[title={What you've built}]
One more word for getting around.
\end{grammarlens}

\subsection*{Your turn}
\begin{itemize}
\raggedright
  \item \emph{Say it:} \emph{naṅkūram}
  \item \emph{Say it:} \emph{naṅkūram}, once more
  \item \emph{Say it:} \emph{kappal tuṟaimukattukku veḷiyē naṅkūram pōṭṭatu.} — The ship dropped anchor outside the harbour.
  \item \emph{Recall:} say the Tamil for a net, then the Tamil for a sailor, then say \emph{naṅkūram} again
\end{itemize}

\begin{tcolorbox}[breakable,colback=teal!4,colframe=teal!35!black,title={Before you move on}]
What does \ta{நங்கூரம்} mean? (\textquotedblleft{}An anchor\textquotedblright{}.) Say it once more.
\end{tcolorbox}
\section[\texorpdfstring{\ta{தெப்பம்} (teppam)}{தெப்பம் (teppam)}]{\ta{தெப்பம்} (teppam) — a raft; a float}

\label{lesson:TA-C417-teppam}



\begin{quote}
Before the new one: say the Tamil for a sailor, then the Tamil for an anchor.
\end{quote}

\subsection*{You'll want to know: \ta{தெப்பம்}}
\textbf{\ta{தெப்பம்}} — \emph{teppam} — \textquotedblleft{}a raft; a float\textquotedblright{}.

\begin{quote}
\ta{சிறுவர்கள்} \ta{மூங்கிலால்} \ta{ஒரு} \ta{தெப்பம்} \ta{கட்டினார்கள்}.
\end{quote}

\begin{quote}
\emph{ciṟuvarkaḷ mūṅkilāl oru teppam kaṭṭiṉārkaḷ.}
\end{quote}

\begin{quote}
The boys built a raft of bamboo.
\end{quote}

Made of bamboo or logs tied together.

\begin{cousinweb}
\textbf{\ta{தெப்பம்}} (\emph{teppam}) is native Dravidian (DEDR 3414), \textquotedblleft{}a raft, a float\textquotedblright{}.
\end{cousinweb}

\begin{grammarlens}[title={What you've built}]
One more word for getting around.
\end{grammarlens}

\subsection*{Your turn}
\begin{itemize}
\raggedright
  \item \emph{Say it:} \emph{teppam}
  \item \emph{Say it:} \emph{teppam}, once more
  \item \emph{Say it:} \emph{ciṟuvarkaḷ mūṅkilāl oru teppam kaṭṭiṉārkaḷ.} — The boys built a raft of bamboo.
  \item \emph{Recall:} say the Tamil for a sailor, then the Tamil for an anchor, then say \emph{teppam} again
\end{itemize}

\begin{tcolorbox}[breakable,colback=teal!4,colframe=teal!35!black,title={Before you move on}]
What does \ta{தெப்பம்} mean? (\textquotedblleft{}A raft; a float\textquotedblright{}.) Say it once more.
\end{tcolorbox}
\section[\texorpdfstring{\ta{தத்தளி} (tattaḷi)}{தத்தளி (tattaḷi)}]{\ta{தத்தளி} (tattaḷi) — to flounder; to be tossed about}

\label{lesson:TA-C417-tattali}



\begin{quote}
Before the new one: say the Tamil for an anchor, then the Tamil for a raft.
\end{quote}

\subsection*{You'll want to know: \ta{தத்தளி}}
\textbf{\ta{தத்தளி}} — \emph{tattaḷi} — \textquotedblleft{}to flounder; to be tossed about\textquotedblright{}.

\begin{quote}
\ta{புயலில்} \ta{சிறிய} \ta{படகு} \ta{தத்தளித்தது}.
\end{quote}

\begin{quote}
\emph{puyalil ciṟiya paṭaku tattaḷittatu.}
\end{quote}

\begin{quote}
The small boat was tossed about in the storm.
\end{quote}

Past (I …): \ta{தத்தளித்தேன்}, \emph{tattaḷittēṉ}.

Also of people in trouble: \ta{பணம்} \ta{இல்லாமல்} \ta{தத்தளித்தேன்}.

\begin{cousinweb}
\textbf{\ta{தத்தளி}} (\emph{tattaḷi}) is native Dravidian (DEDR 3061), \textquotedblleft{}to be greatly agitated, to flounder\textquotedblright{}.
\end{cousinweb}

\begin{grammarlens}[title={What you've built}]
One more word for getting around.
\end{grammarlens}

\subsection*{Your turn}
\begin{itemize}
\raggedright
  \item \emph{Say it:} \emph{tattaḷi}
  \item \emph{Say it:} \emph{tattaḷi}, once more
  \item \emph{Say it:} \emph{puyalil ciṟiya paṭaku tattaḷittatu.} — The small boat was tossed about in the storm.
  \item \emph{Recall:} say the Tamil for an anchor, then the Tamil for a raft, then say \emph{tattaḷi} again
\end{itemize}

\begin{tcolorbox}[breakable,colback=teal!4,colframe=teal!35!black,title={Before you move on}]
What does \ta{தத்தளி} mean? (\textquotedblleft{}To flounder; to be tossed about\textquotedblright{}.) Say it once more.
\end{tcolorbox}
\section[\texorpdfstring{\ta{கரை} (karai)}{கரை (karai)}]{\ta{கரை} (karai) — a shore; a bank}

\label{lesson:TA-C417-karai}



\begin{quote}
Before the new one: say the Tamil for a raft, then the Tamil for to flounder.
\end{quote}

\subsection*{You'll want to know: \ta{கரை}}
\textbf{\ta{கரை}} — \emph{karai} — \textquotedblleft{}a shore; a bank\textquotedblright{}.

\begin{quote}
\ta{மாலுமிகள்} \ta{மாலையில்} \ta{கரைக்குத்} \ta{திரும்பினார்கள்}.
\end{quote}

\begin{quote}
\emph{mālumikaḷ mālaiyil karaikkut tirumpiṉārkaḷ.}
\end{quote}

\begin{quote}
The sailors came back to shore in the evening.
\end{quote}

\ta{கடற்கரை} is the beach; \ta{ஆற்றங்கரை}, a riverbank.

\begin{cousinweb}
\textbf{\ta{கரை}} (\emph{karai}) is native Dravidian (DEDR 1293), \textquotedblleft{}a shore, a bank, a ridge, a border\textquotedblright{}.
\end{cousinweb}

\begin{grammarlens}[title={What you've built}]
One more word for getting around.
\end{grammarlens}

\subsection*{Your turn}
\begin{itemize}
\raggedright
  \item \emph{Say it:} \emph{karai}
  \item \emph{Say it:} \emph{karai}, once more
  \item \emph{Say it:} \emph{mālumikaḷ mālaiyil karaikkut tirumpiṉārkaḷ.} — The sailors came back to shore in the evening.
  \item \emph{Recall:} say the Tamil for a raft, then the Tamil for to flounder, then say \emph{karai} again
\end{itemize}

\begin{tcolorbox}[breakable,colback=teal!4,colframe=teal!35!black,title={Before you move on}]
What does \ta{கரை} mean? (\textquotedblleft{}A shore; a bank\textquotedblright{}.) Say it once more.
\end{tcolorbox}
